% Separate post-CP21 candidate; numerical least-index claim needs full audits.
\subsection{Second-level point fields close the reducible branch}

Retain the exact mixed-characteristic, residue, determinant and
valuation conventions of B60.

\begin{lemma}[The elementary-abelian point field]
\label{lem:B-RD61-point-field}
\status{Candidate; complete proof written, independent audits pending.}
Let $R$ contain the full $p$-torsion of a height-two group $M$, and
let $P$ have order $p^2$. Then $E=R(P)$ is Galois over $R$, with
elementary abelian group of order at most $p^2$. If its degree is
$p^2$, the map
\[
       \operatorname{Gal}(E/R)\longrightarrow M[p],
                       \qquad g\longmapsto gP-P        \tag{B-RD61.1}
\]
is an isomorphism of $\mathbf F_p$-spaces. An isomorphism of the
actual generic $p^2$-torsion modules identifies the point fields
and this map for corresponding points.
\end{lemma}
\begin{proof}
Over $R$, the image on $M[p^2]\simeq(\mathbf Z/p^2)^2$ lies
in $1+pM_2(\mathbf F_p)$, an elementary abelian group.
Consequently every subfield of the full $p^2$-division field is
Galois over $R$. All conjugates of $P$ lie in the fibre of $[p]$
over the rational point $pP$, a set of $p^2$ points. This bounds
the degree of $R(P)$. The displayed map is an injective homomorphism
because all translations are $R$-rational; at degree $p^2$ it is
surjective. Finally an equivariant torsion-module isomorphism gives
exactly the same stabilizer for corresponding points and translates,
which proves the last assertion. The coordinate field of a point
in the good formal group equals its field of definition: an
integral parameter is a formal coordinate, so fixing it fixes the
point, and conversely.
\end{proof}

\begin{proposition}[No remaining wild reducible comparison]
\label{prop:B-RD61-wild}
\status{Candidate; complete proof written, independent audits pending.}
Two normalized potentially supersingular height-two groups with
isomorphic $p^2$-torsion and reducible residual inertia of nontrivial
wild image have the same special character.
\end{proposition}
\begin{proof}
By Proposition~\ref{prop:B-RD60-residual-break}, only the case
$p\nmid B_1$, $p\mid B_2$ and $B_1>B_2$ could have different
characters. Identify their torsion modules, and let $R/L$ be their
common full $p$-division field, cyclic of degree $p$.
Their unique high lines coincide; call the common line $H$.
Choose corresponding points $P_i$ of order $p^2$ with
$pP_i\notin H$. Formula (B-RD60.5) gives
\[
                 v_R(T_i(P_i))=B_i/p^2.
\]
On the first side this fraction has denominator $p^2$. The
ramification index of its point field is therefore divisible by
$p^2$. Lemma~\ref{lem:B-RD61-point-field} bounds its degree by
$p^2$, so the common point field $E=R(P_i)$ is Galois of degree
exactly $p^2$. In particular
\[
       v_E(T_i(P_i))=B_i,
       \qquad [E:L]=p^3.                              \tag{B-RD61.2}
\]
Every $p$-torsion translation occurs in (B-RD61.1).
Choose $g$ whose nonzero translation lies outside $H$.
Its coordinate difference in model $i$ has valuation
\[
       [E:L]\,\frac{B_i}{p}=p^2 B_i.                 \tag{B-RD61.3}
\]
Since $p\nmid B_1$, Lemma~\ref{lem:B-RD60-break} identifies the
actual lower break of this particular $g$ on $E/R$ as
\[
                          b_E(g)=NB_1.
\]
Applying just the inequality of that lemma to the second model
then gives
\[
                         p^2B_2\ge B_2+NB_1,
\]
which contradicts $B_1>B_2$. This excludes the final wild case.
The argument uses an actual common point field, not an abstract
field with the same order or a substituted integral lattice.
\end{proof}

\begin{proposition}[No remaining tame reducible comparison]
\label{prop:B-RD61-tame}
\status{Candidate; complete proof written, independent audits pending.}
Two normalized potentially supersingular height-two groups with
isomorphic $p^2$-torsion and tame reducible residual inertia have
the same special character.
\end{proposition}
\begin{proof}
Here the common $p$-division field over $L$ is $L$ itself.
Write $B_i=pb_i$, with $1\le b_i\le p-2$, by B60.
Their high lines $H_1,H_2$ are not initially identified.
Choose a common nonzero residual vector outside their union,
and lift it to corresponding points $P_i$ of order $p^2$.
This is possible because the union of two lines is a proper
subset of $\mathbf F_p^2$. The low primitive valuations are
\[
                          v_L(T_i(P_i))=b_i/p^2.
\]
As $p\nmid b_i$, their common point field $E=L(P_i)$ has degree
$p^2$ by Lemma~\ref{lem:B-RD61-point-field} and the valuation
denominator. Hence $v_E(T_i(P_i))=b_i$ and all translations in
$M_i[p]$ occur.

Choose a translation outside both high lines. Its coordinate
difference in model $i$ has valuation $p^2b_i$.
The equality part of Lemma~\ref{lem:B-RD60-break} therefore gives
one and the same lower break
\[
                       b_E(g)=Nb_1=Nb_2.
\]
Thus $b_1=b_2=b$. If the high lines were different, choose a
nonzero translation lying in $H_1$ and outside $H_2$.
The corresponding difference valuations would be, respectively,
\[
                   p^2(N-pb)\quad\hbox{and}\quad p^2b.
\]
Both coordinates of $P_i$ have the same prime-to-$p$ valuation
$b$, so their differences under this same $g$ must have the same
valuation by the equality part of the lower-break lemma. That
would force $N-pb=b$, or $b=p-1$, outside its allowed range.
Consequently $H_1=H_2$. Their common character
$\zeta^{c_i-pb}$ in (B-RD60.4) now gives $c_1=c_2$ modulo $N$,
as required.
\end{proof}

\begin{theorem}[The least index at minimal cyclotomic ramification]
\label{thm:B-RD61}
\status{Candidate; complete proof written, independent audits pending.}
For every prime $p\ge5$, under the manuscript's monodromy condition,
height two and $e_K=p-1$, with $\mu_p\subset K$ and perfect residue
field, the exact least index in B.4 is
\[
                         \nu_{\min}(p,p-1,2)=1.         \tag{B-RD61.4}
\]
Thus $p^2$-torsion always determines the requested special
prime-to-$p$ inertia character, while $p$-torsion need not.
In particular $\nu_{\min}(5,4,2)=1$.
\end{theorem}
\begin{proof}
The lower bound is the actual first-level counterexample of B18,
as used in Theorem~\ref{thm:B-RD49}. For the upper bound pass to
the completed maximal unramified extension, without changing
absolute ramification or the character question on inertia.
The dimension and mixed-case reductions of B49 leave only two
whole potentially supersingular groups with the same finite
determinant. Proposition~\ref{prop:B-RD50-square} handles its
square case. In the nonsquare case the common finite height-one
twist normalizes the determinant exactly as in B53 and B60.
The common residual module is either irreducible over $\mathbf F_p$,
handled by Theorem~\ref{thm:B-RD59}, or reducible. Its wild image
is then trivial or cyclic of order $p$; the two preceding
propositions handle these cases. The special characters agree
in every case, proving the upper bound and equality.
\end{proof}

\begin{remark}[Scope of the new equality]
\label{rem:B-RD61-scope}
\status{Candidate scope statement.}
The equality treats minimal cyclotomic ramification and height two
for $p\ge5$, including non-CM deformations and reducible residual
representations. It does not determine the general B.4 least index
at other ramification indices or heights. All point-field degrees
used above follow from the actual torsion action and reduced
valuation denominators. No existence conclusion is drawn from
an enumeration of possible Newton slopes.
The exact finite supplement
\path{code/reduction_descent/verify_tame_reducible_recovery.py}
checks 74,828 Newton/tame-character records for every prime from
5 through 101, 73,696 prime-to-$p$ wild records, 1,132 divisible
wild comparisons, 36,848 mixed point-field inequalities, 1,698
tame residual pairs (including 566 distinct-type exclusions), and
5,184 exact integer-binomial bounds. The author run passed in
1.8442117410013452 seconds. Its separate exploratory $p=5$ program
and full output are retained: the 20 Newton records give eight
necessary wild rows and two necessary tame rows before the
point-field contradiction. Neither finite list is used as a
substitute for the general proof.
\end{remark}
